%% Annexe A : Glossaire
\chapter{Glossaire}
\label{ann:glossaire}

Ce glossaire reprend les notions définies dans les chapitres. Le renvoi
indique le chapitre où la notion est présentée.

\begin{description}
    \item[Admissibilité] Aptitude d'une preuve à être reçue par un tiers, selon la manière dont elle a été recueillie et documentée (Ch.~\ref{chap:preuve}).
    \item[Analyse forensique] Recherche méthodique, documentée et opposable des traces d'un événement (Ch.~\ref{chap:fondements}).
    \item[ANTIC] Agence Nationale des Technologies de l'Information et de la Communication, autorité camerounaise de la sécurité électronique (Ch.~\ref{chap:droit}).
    \item[Canal de pilotage] Connexions régulières de faible volume par lesquelles un attaquant commande un code malveillant (Ch.~\ref{chap:reseau}).
    \item[CERT] Appellation protégée d'un CSIRT~; l'ANTIC opère le CERT national (Ch.~\ref{chap:csirt}).
    \item[Chaîne d'attaque] Description d'une intrusion en étapes, de la reconnaissance à l'action sur objectif (Ch.~\ref{chap:renseignement}).
    \item[Chaîne de custody] Continuité documentée de la détention d'une preuve (Ch.~\ref{chap:fondements}).
    \item[Chasse] Recherche proactive de ce qui n'a pas déclenché d'alerte (Ch.~\ref{chap:soc}).
    \item[Condensat] Empreinte numérique qui prouve l'intégrité d'une copie (Ch.~\ref{chap:fondements}).
    \item[Corréler] Rapprocher automatiquement des événements de sources différentes pour faire émerger une alerte (Ch.~\ref{chap:siem}).
    \item[CSIRT] Équipe de réponse aux incidents de sécurité d'un périmètre (Ch.~\ref{chap:csirt}).
    \item[Culture juste] Culture où l'on signale les erreurs sans crainte de sanction injuste (Ch.~\ref{chap:retex}).
    \item[Délai de présence] Durée entre la première compromission et sa détection (Ch.~\ref{chap:pourquoi}).
    \item[Exfiltration] Transfert sortant anormal de données vers une destination contrôlée par l'attaquant (Ch.~\ref{chap:reseau}).
    \item[Faux négatif] Comportement visé qu'une règle laisse passer (Ch.~\ref{chap:regles}).
    \item[Faux positif] Alerte sans incident réel (Ch.~\ref{chap:regles}).
    \item[Indicateur de compromission] Élément observable (adresse, condensat, domaine) trahissant une attaque (Ch.~\ref{chap:renseignement}).
    \item[Ingénierie de détection] Écriture et maintenance des règles de détection (Ch.~\ref{chap:regles}).
    \item[ISO/IEC 27035] Norme de gestion des incidents de sécurité, cycle en cinq phases (Ch.~\ref{chap:normes}).
    \item[ISO/IEC 27037] Lignes directrices de collecte et de préservation de la preuve numérique (Ch.~\ref{chap:normes}).
    \item[MITRE ATT\&CK] Base de connaissances des tactiques et techniques des attaquants (Ch.~\ref{chap:renseignement}).
    \item[Modèle du Diamant] Cadre d'analyse d'une intrusion~: adversaire, capacités, infrastructure, victime (Ch.~\ref{chap:renseignement}).
    \item[NIST SP 800-86] Guide d'intégration de la forensique à la réponse~: collecte, examen, analyse, rapport (Ch.~\ref{chap:normes}).
    \item[Ordre de volatilité] Ordre d'acquisition des preuves, du plus fugace au plus stable (Ch.~\ref{chap:fondements}).
    \item[RFC 3227] Lignes directrices de collecte et d'archivage des preuves (Ch.~\ref{chap:normes}).
    \item[SIEM] Plateforme de collecte, corrélation et alerte sur les journaux (Ch.~\ref{chap:siem}).
    \item[SOC] Centre d'opérations de sécurité~: détecter, analyser, traiter en continu (Ch.~\ref{chap:soc}).
    \item[STIX / TAXII] Formats d'échange du renseignement sur la menace (Ch.~\ref{chap:renseignement}).
    \item[Supervision de sécurité] Collecte et analyse continues des traces pour détecter les attaques (Ch.~\ref{chap:pourquoi}).
    \item[Trilemme CRO] Confidentialité, Fiabilité, Opposabilité~: les trois dimensions de toute trace et de tout acte d'enquête (Ch.~\ref{chap:prologue}).
    \item[Verdict à trois valeurs] OUI (démontré), NON (réfuté), NEK (indéterminé) (Ch.~\ref{chap:mesurer}).
    \item[Vrai positif] Alerte sur un comportement réellement visé (Ch.~\ref{chap:regles}).
\end{description}
